% Luc Destombe --- licence CC BY-NC-SA 4.0
% https://creativecommons.org/licenses/by-nc-sa/4.0/deed.fr
% Fichier autonome : compiler avec pdflatex, deux fois (signets, nombre de pages).
\def\lyceeversion{1.0}
\documentclass[a4paper,10pt]{article}
\makeatletter
\newif\iflycee@niveaux
\newif\iflycee@palatino
\lycee@palatinotrue

\RequirePackage[utf8]{inputenc}
\RequirePackage[T1]{fontenc}
\RequirePackage[french]{babel}
\RequirePackage{etoolbox}

\newcommand{\ShorthandsOff}{\@ifundefined{shorthandoff}{}{\shorthandoff{;:!?}}}
\newcommand{\ShorthandsOn}{\@ifundefined{shorthandon}{}{\shorthandon{;:!?}}}
\ShorthandsOff
\AtBeginDocument{\ShorthandsOn}
\AtBeginEnvironment{tikzpicture}{\ShorthandsOff}

\RequirePackage{geometry}
\RequirePackage{fancyhdr}
\RequirePackage{lastpage}
\RequirePackage{multicol}
\RequirePackage{array}
\RequirePackage{enumitem}
\RequirePackage{xcolor}
\RequirePackage{needspace}

\RequirePackage{amsmath,amssymb}
\RequirePackage{mathpazo}
\renewcommand{\ttdefault}{lmtt}
\DeclareMathSizes{10.95}{10.95}{8}{5.5}
\begingroup\catcode`\_=\active \catcode`\^=\active
\gdef_#1{\sb{\scriptscriptstyle #1}}
\gdef^#1{\sp{\scriptscriptstyle\lycee@moinsepais #1}}
\endgroup
\newcommand\lycee@moins{\mathbin{%
  \setbox\z@\hbox{$\m@th\scriptscriptstyle\lycee@moinsorig$}%
  \dimen@=.55\wd\z@ \advance\dimen@-.5pt
  \hbox to.75\wd\z@{\hss$\m@th\scriptscriptstyle\vcenter{\hbox{%
    \tikz\draw[line width=.5pt,line cap=round](0,0)--(\the\dimen@,0);}}$\hss}}}
\begingroup\lccode`\~=`\- \lowercase{\endgroup\let~}\lycee@moins
\newcommand\lycee@moinsepais{\mathcode`\-="8000 }
\AtBeginDocument{\mathchardef\lycee@moinsorig=\mathcode`\-\relax}
\AtBeginDocument{\catcode`\_=12 \mathcode`\_="8000
  \catcode`\^=12 \mathcode`\^="8000 }
\newcommand\lycee@exposants{%
  \fontdimen13\textfont2=.48\fontdimen6\textfont2
  \fontdimen14\textfont2=.48\fontdimen6\textfont2
  \fontdimen15\textfont2=.48\fontdimen6\textfont2
  \fontdimen13\scriptfont2=.48\fontdimen6\scriptfont2
  \fontdimen14\scriptfont2=.48\fontdimen6\scriptfont2
  \fontdimen15\scriptfont2=.48\fontdimen6\scriptfont2}
\every@math@size\expandafter{\the\every@math@size\lycee@exposants}
\newlength{\retraitcalculs}
\setlength{\retraitcalculs}{2em}

\RequirePackage{tikz}
\usetikzlibrary{arrows.meta,calc,babel,shapes.misc}
\RequirePackage[most]{tcolorbox}
\tcbuselibrary{skins,breakable}

\definecolor{coulDef}{HTML}{1F4E79}
\definecolor{coulProp}{HTML}{7030A0}
\definecolor{coulRem}{HTML}{C55A11}
\definecolor{coulEx}{HTML}{2E7D32}
\definecolor{coulExo}{HTML}{B03A2E}
\definecolor{coulPart}{HTML}{1F4E79}
\definecolor{coulMeth}{HTML}{1B6E4F}
\definecolor{coulCourbe}{HTML}{0B5ED7}
\definecolor{coulCourbeB}{HTML}{C00000}
\definecolor{coulCourbeC}{HTML}{2E7D32}
\definecolor{coulGrille}{HTML}{8C8C8C}
\definecolor{coulRep}{HTML}{1B6E4F}
\definecolor{coulBar}{HTML}{2E7D32}

\newcommand{\R}{\mathbb{R}}
\newcommand{\Rs}{\mathbb{R}^{*}}
\renewcommand{\le}{\leqslant}
\renewcommand{\ge}{\geqslant}
\renewcommand{\approx}{\simeq}
\newcommand{\vect}[1]{\overrightarrow{#1}}
\newcommand{\norme}[1]{\left\lVert #1 \right\rVert}
\newcommand{\intervalle}[2]{\left[#1\,;#2\right]}
\RequirePackage{cancel}
\renewcommand{\CancelColor}{\color{coulExo}}
\newcommand{\fois}[1]{\times{\color{coulExo}#1}}
\newcommand{\barre}[1]{\cancel{#1}}
\newlength{\lycee@eqh}\newlength{\lycee@eqd}
\newsavebox{\lycee@eqa}\newsavebox{\lycee@eqb}\newsavebox{\lycee@eqc}
\newcommand{\lycee@eqcol}[1]{\hbox{$\renewcommand{\arraystretch}{1.35}%
  \begin{array}[t]{@{}>{\displaystyle}r@{}>{\displaystyle{}}l@{}}#1\end{array}$}}
\newcommand{\lycee@eqmesure}[1]{%
  \ifdim\ht#1>\lycee@eqh\setlength{\lycee@eqh}{\ht#1}\fi
  \ifdim\dp#1>\lycee@eqd\setlength{\lycee@eqd}{\dp#1}\fi}
\newcommand{\lycee@eqtrait}{\hspace{1.2em}\tikz[baseline=0pt]\draw[dash pattern=on 2pt off 2pt,
  line width=0.5pt] (0,\lycee@eqh) -- (0,-\lycee@eqd);\hspace{1.2em}}
\newcommand{\eqproduit}[2]{\par\smallskip\noindent C'est une équation produit :
  $A\times B=0 \iff A=0$ ou $B=0$.\par\nopagebreak\smallskip
  \begingroup\setlength{\lycee@eqh}{0pt}\setlength{\lycee@eqd}{0pt}%
  \setbox\lycee@eqa=\lycee@eqcol{#1}\setbox\lycee@eqb=\lycee@eqcol{#2}%
  \lycee@eqmesure\lycee@eqa\lycee@eqmesure\lycee@eqb%
  \noindent\hspace*{\retraitcalculs}\usebox\lycee@eqa\lycee@eqtrait\usebox\lycee@eqb\endgroup\par\smallskip}
\newcommand{\eqproduittrois}[3]{\par\smallskip\noindent C'est une équation produit :
  $A\times B\times C=0 \iff A=0$ ou $B=0$ ou $C=0$.\par\nopagebreak\smallskip
  \begingroup\setlength{\lycee@eqh}{0pt}\setlength{\lycee@eqd}{0pt}%
  \setbox\lycee@eqa=\lycee@eqcol{#1}\setbox\lycee@eqb=\lycee@eqcol{#2}\setbox\lycee@eqc=\lycee@eqcol{#3}%
  \lycee@eqmesure\lycee@eqa\lycee@eqmesure\lycee@eqb\lycee@eqmesure\lycee@eqc%
  \noindent\hspace*{\retraitcalculs}\usebox\lycee@eqa\lycee@eqtrait\usebox\lycee@eqb\lycee@eqtrait\usebox\lycee@eqc\endgroup
  \par\smallskip}

\tikzset{grille/.style={coulGrille,line width=0.4pt}}
\newcommand{\quadrillage}[4]{\draw[grille,step=1] (#1,#2) grid (#3,#4);}
\tikzset{
  croix/.style={cross out,draw,line width=0.6pt,inner sep=0pt,minimum size=4pt},
  croixdroite/.style={croix,rotate=45,line width=0.8pt},
}
\newcommand{\pt}[3]{\path (#1) node[croix]{} node[#2,font=\small,inner sep=2.5pt]{$#3$};}

\newcommand{\lycee@repere}[4]{%
  \draw[grille,step=1] (#1,#2) grid (#3,#4);
  \draw[-{Stealth[length=2mm]},thick] (#1,0)--({#3+0.4},0) node[below left=-1pt]{$x$};
  \draw[-{Stealth[length=2mm]},thick] (0,#2)--(0,{#4+0.4}) node[below left=-1pt]{$y$};
  \node[below left,font=\scriptsize,inner sep=1.5pt] at (0,0) {$0$};}
\newcommand{\lycee@gradx}[1]{\foreach \k/\l in {#1}{%
  \draw (\k,0.07)--(\k,-0.07) node[below,font=\scriptsize,inner sep=1.5pt]{$\l$};}}
\newcommand{\lycee@grady}[1]{\foreach \k/\l in {#1}{%
  \draw (0.07,\k)--(-0.07,\k) node[left,font=\scriptsize,inner sep=1.5pt]{$\l$};}}
\newcommand{\lycee@intff}[2]{\left[#1\,;#2\right]}
\newcommand{\lycee@intfo}[2]{\left[#1\,;#2\right[}
\newcommand{\lycee@intof}[2]{\left]#1\,;#2\right]}
\newcommand{\lycee@intoo}[2]{\left]#1\,;#2\right[}
\newcommand{\lycee@ens}[1]{\left\{#1\right\}}
\AtBeginDocument{%
  \providecommand{\repere}{\lycee@repere}%
  \providecommand{\gradx}{\lycee@gradx}%
  \providecommand{\grady}{\lycee@grady}%
  \providecommand{\Cf}{\mathcal{C}\sb{\scriptscriptstyle f}}%
  \providecommand{\Cg}{\mathcal{C}\sb{\scriptscriptstyle g}}%
  \providecommand{\intff}{\lycee@intff}%
  \providecommand{\intfo}{\lycee@intfo}%
  \providecommand{\intof}{\lycee@intof}%
  \providecommand{\intoo}{\lycee@intoo}%
  \providecommand{\ens}{\lycee@ens}}

\newlist{questions}{enumerate}{3}
\setlist[questions,1]{label=\arabic*\textdegree),leftmargin=*,itemsep=2pt,topsep=3pt}
\setlist[questions,2]{label=\alph*),leftmargin=*,itemsep=1pt,topsep=2pt}
\setlist[questions,3]{label=\roman*),leftmargin=*,itemsep=1pt,topsep=1pt}
\setlist[itemize]{label=\textbullet,leftmargin=*,itemsep=2pt,topsep=3pt}

\newcommand{\besoinplace}[1]{\par\penalty\@M\begingroup
  \dimen@=#1\relax\dimen@ii=\pagegoal\advance\dimen@ii-\pagetotal
  \ifdim\dimen@>\dimen@ii\ifdim\dimen@ii>\z@\vfil\fi\break\fi\endgroup}

\newcommand{\lycee@pied}{\thepage/\pageref{LastPage}}
\newcommand{\entete}[2]{%
  \pagestyle{fancy}\fancyhf{}%
  \fancyhead[L]{\footnotesize\color{coulPart}\textbf{#1}}%
  \fancyhead[R]{\footnotesize\color{coulPart}\textbf{#2}}%
  \fancyfoot[C]{\footnotesize\lycee@pied}%
  \renewcommand{\headrulewidth}{0.4pt}%
  \renewcommand{\headrule}{\hbox to\headwidth{%
    \color{coulPart!40}\leaders\hrule height \headrulewidth\hfill}}}

\RequirePackage{ccicons}
\newcommand{\lycee@licence}{{\scriptsize\color{gray}%
  \copyright\ Luc Destombe --- \ccbyncsa\ CC BY-NC-SA 4.0}}
\newif\iflycee@licence \lycee@licencetrue
\newcommand{\sanslicence}{\lycee@licencefalse}
\AtBeginDocument{\iflycee@licence\AddToHookNext{shipout/foreground}{%
  \put(0,0){\rlap{\kern\dimexpr1in+\hoffset+\oddsidemargin+\textwidth\relax
    \llap{\raisebox{-\dimexpr1in+\voffset+\topmargin+\headheight+\headsep
      +\textheight+\footskip\relax}{\lycee@licence}}}}}\fi}

\newcommand{\lycee@site}{faire-des-maths.fr}
\newcommand{\lycee@version}{{\scriptsize\color{gray}\lycee@site
  \ifdefined\lyceeversion\ --- Version~\lyceeversion\fi}}
\AtBeginDocument{\iflycee@licence\AddToHookNext{shipout/foreground}{%
  \put(0,0){\rlap{\kern\dimexpr1in+\hoffset+\oddsidemargin\relax
    \raisebox{-\dimexpr1in+\voffset+\topmargin+\headheight+\headsep
      +\textheight+\footskip\relax}{\lycee@version}}}}\fi}

\newcommand{\formulesserrees}{%
  \setlength{\abovedisplayskip}{3pt}\setlength{\belowdisplayskip}{3pt}%
  \setlength{\abovedisplayshortskip}{2pt}\setlength{\belowdisplayshortskip}{3pt}}

\setlength{\parindent}{0pt}

\RequirePackage[implicit=false,hidelinks,pdfencoding=auto,psdextra,bookmarksopen,
  bookmarksopenlevel=1,pdfcreator={},pdfproducer={}]{hyperref}
\RequirePackage{bookmark}
\hypersetup{pdfauthor={Luc Destombe},pdfkeywords={CC BY-NC-SA 4.0}}
\newcounter{lycee@signet}
\newcommand{\signet}[3][]{\stepcounter{lycee@signet}%
  \edef\lycee@dest{\if\relax\detokenize{#1}\relax
    signet.\arabic{lycee@signet}\else#1\fi}%
  \hypertarget{\lycee@dest}{}\bookmark[level=#2,dest=\lycee@dest]{#3}}
\pdfstringdefDisableCommands{%
  \def\R{R}\def\vect#1{#1}\def\Cf{Cf}\def\Cg{Cg}\def\fois#1{×#1}%
  \def\barre#1{#1}\def\intervalle#1#2{[#1;#2]}\def\intff#1#2{[#1;#2]}%
  \def\textsuperscript#1{#1}\def\dfrac#1#2{#1/#2}\def\frac#1#2{#1/#2}%
  \def\vec#1{#1⃗}}%
\RequirePackage[official]{eurosym}

\geometry{top=0.8cm,bottom=1.3cm,left=1cm,right=1cm,footskip=0.6cm}

\pagestyle{fancy}
\fancyhf{}
\renewcommand{\headrulewidth}{0pt}
\fancyfoot[C]{\footnotesize Page \thepage{} / \pageref{LastPage}}

\newcommand{\titrefiche}[2]{%
  \begin{center}
    \Large\textbf{#1}\\[2pt]
    \large\textbf{#2}\\[2pt]
  \end{center}
  \vspace{0.10cm}}

\setlength{\columnseprule}{0.4pt}
\setlength{\columnsep}{15pt}
\renewcommand{\columnseprulecolor}{\color{black!45}}
\setlength{\multicolsep}{2pt}

\newcounter{partiefiche}
\renewcommand{\thepartiefiche}{\Roman{partiefiche}}
\newif\ifapresbandeau
\newcommand{\lycee@bandeau}[1]{%
  \par\penalty-600
  \vspace{0.08cm}%
  \noindent\colorbox{coulPart}{%
    \parbox{\dimexpr\linewidth-2\fboxsep\relax}{%
      \color{white}\bfseries\raggedright\strut#1\strut}}%
  \nopagebreak\par\nopagebreak
  \vspace{0.06cm}\nopagebreak
  \apresbandeautrue}
\newcommand{\partiefiche}[1]{%
  \refstepcounter{partiefiche}%
  \lycee@bandeau{\signet{1}{\thepartiefiche{} --- #1}\thepartiefiche{} --- #1}}
\newcommand{\partiefichesansnum}[1]{\lycee@bandeau{\signet{1}{#1}#1}}

\newcounter{exo}
\newlength{\espaceexo}\setlength{\espaceexo}{7pt}
\newlength{\espacetitre}\setlength{\espacetitre}{2pt}
\newcommand{\exo}[1][\relax]{%
  \par
  \ifapresbandeau\apresbandeaufalse\else\penalty-150\addvspace{\espaceexo}\fi
  \refstepcounter{exo}%
  \noindent\signet[exercice-\arabic{exo}]{2}{Exercice \arabic{exo}}%
  \textbf{\color{coulExo}\underline{Exercice \theexo}}%
  \ifx\relax#1\relax\else\textbf{ : #1}\fi
  \par\nopagebreak\vspace{\espacetitre}\nopagebreak
  \ignorespaces}

\setlist[enumerate]{nosep,leftmargin=*,itemsep=2pt,topsep=2pt}
\setlist[itemize]{nosep,leftmargin=*,topsep=2pt}
\newcommand{\choix}[3]{\par\hspace*{1em}%
  \makebox[0.3\linewidth][l]{a) #1}\makebox[0.3\linewidth][l]{b) #2}c) #3}
\newcommand{\deux}[2]{\par\hspace*{0.6em}%
  \makebox[0.48\linewidth][l]{#1}#2}
\newcommand{\trois}[3]{\par\hspace*{0.6em}%
  \makebox[0.32\linewidth][l]{#1}\makebox[0.32\linewidth][l]{#2}#3}

\renewenvironment{center}{\par\addvspace{3pt}\centering}{\par\addvspace{3pt}}
\formulesserrees
\AtBeginDocument{\formulesserrees}
\clubpenalty=10000
\widowpenalty=10000
\displaywidowpenalty=10000
\brokenpenalty=10000
\setlength{\parskip}{0pt}

\RequirePackage{ragged2e}
\setlength{\RaggedRightRightskip}{0pt plus 1fil}
\AtBeginDocument{\RaggedRight\binoppenalty=10000 \relpenalty=10000 }
\g@addto@macro\@parboxrestore{\RaggedRight}
\makeatother
\geometry{bottom=1cm}

\setlength{\espaceexo}{4pt}
\setlist[enumerate]{nosep,leftmargin=*,itemsep=1pt,topsep=1pt}

\newenvironment{tableaustat}{\par\smallskip\centering\small
  \setlength{\tabcolsep}{3.5pt}\renewcommand{\arraystretch}{1.3}}{\par\smallskip}
\tikzset{barre/.style={fill=coulCourbe!20,draw=coulCourbe,line width=0.8pt},
  secteur/.style={draw=black,line width=0.6pt}}

\begin{document}

\titrefiche{Seconde --- Statistiques}{Fiche d'exercices du chapitre}

\begin{multicols}{2}\raggedcolumns

\partiefiche{Effectifs, fréquences, effectifs cumulés, classes}

\exo[Activités hors du lycée]
Dans une classe, $26$ élèves ont répondu à la question « À combien d'activités
(club, association, cours de musique\ldots) êtes-vous inscrit ? ». Un élève a répondu
« ne se prononce pas » ; voici les $25$ autres réponses.
\begin{center}
$1$ ; $0$ ; $2$ ; $1$ ; $1$ ; $3$ ; $0$ ; $2$ ; $1$ ; $0$ ; $1$ ; $2$ ; $1$ ; $3$ ;
$0$ ; $1$ ; $2$ ; $0$ ; $1$ ; $1$ ; $2$ ; $3$ ; $0$ ; $2$ ; $1$
\end{center}
\begin{enumerate}[label=\alph*)]
  \item Préciser la population, le caractère étudié et son type.
  \item Dresser le tableau des effectifs, des fréquences en pourcentage et des
    effectifs cumulés croissants.
  \item Combien d'élèves sont inscrits à au plus une activité ? Quelle proportion
    des élèves cela représente-t-il ?
  \item Quel pourcentage des élèves sont inscrits à au moins deux activités ?
\end{enumerate}

\exo[Le bon vocabulaire]
\begin{enumerate}[label=\alph*)]
  \item Pour chaque caractère, dire s'il est qualitatif, quantitatif discret ou
    quantitatif continu : la couleur préférée ; le nombre de messages envoyés hier ;
    le temps mis pour courir $100$ m ; la pointure ; la masse d'un cartable ; le
    sport pratiqué.
  \item Une enquête porte sur la durée de sommeil des élèves de seconde d'un lycée.
    Écrire une phrase qui en donne la population, puis une phrase qui en donne le
    caractère et son type.
\end{enumerate}

\exo[Effectifs cumulés d'un quiz]
Voici les notes sur $10$ de $32$ élèves à un quiz.
\begin{tableaustat}
\begin{tabular}{|l|*{8}{c|}}
\hline
Note & $3$ & $4$ & $5$ & $6$ & $7$ & $8$ & $9$ & $10$\\\hline
Effectif & $1$ & $2$ & $4$ & $5$ & $8$ & $6$ & $4$ & $2$\\\hline
Eff. cumulé & $1$ & $3$ & & & & & & \\\hline
\end{tabular}
\end{tableaustat}
\begin{enumerate}[label=\alph*)]
  \item Expliquer pourquoi l'effectif cumulé de la note $5$ vaut $7$.
  \item Compléter la ligne des effectifs cumulés croissants.
  \item Combien d'élèves ont au plus $7$ ? au moins $8$ ?
  \item Quelle proportion des élèves ont au plus $6$ ? Répondre par une phrase.
\end{enumerate}

\exo[Venir au lycée]
Un lycée compte $1~250$ élèves. Le tableau donne la répartition des élèves selon
leur moyen de transport principal.
\begin{tableaustat}
\begin{tabular}{|l|*{5}{c|}}
\hline
Mode & bus & voiture & à pied & vélo & train\\\hline
Fréquence & $44\,\%$ & $24\,\%$ & $18\,\%$ & $8\,\%$ & \\\hline
\end{tabular}
\end{tableaustat}
\begin{enumerate}[label=\alph*)]
  \item Calculer la fréquence des élèves qui viennent en train.
  \item Calculer l'effectif de chaque mode de transport.
  \item Dans un autre lycée, $156$ élèves viennent à vélo, soit $12\,\%$ des élèves.
    Combien ce lycée compte-t-il d'élèves ?
\end{enumerate}

\exo[Distance au lycée]
Voici la distance entre le domicile et le lycée de $20$ élèves, en km.
\begin{center}
$0{,}8$ ; $3{,}5$ ; $1{,}2$ ; $6$ ; $2$ ; $4{,}7$ ; $9{,}1$ ; $1{,}9$ ; $5{,}5$ ; $3$ ;
$0{,}5$ ; $7{,}2$ ; $2{,}6$ ; $4$ ; $1{,}5$ ; $8{,}4$ ; $3{,}8$ ; $6{,}9$ ; $2{,}2$ ; $5$
\end{center}
\begin{enumerate}[label=\alph*)]
  \item Regrouper ces distances en classes d'amplitude $2$ km, de $\intfo{0}{2}$ à
    $\intfo{8}{10}$, dans un tableau d'effectifs.
  \item Calculer la fréquence de chaque classe, en pourcentage.
  \item Quel pourcentage des élèves habitent à moins de $4$ km du lycée ?
  \item Avec le seul tableau, peut-on savoir combien d'élèves habitent à moins de
    $3$ km ? Expliquer.
\end{enumerate}

\partiefiche{Moyenne, moyenne pondérée}

\exo[Trois moyennes]
\begin{enumerate}[label=\alph*)]
  \item Une personne a compté ses pas pendant une semaine, en milliers :
    $8$ ; $11$ ; $6$ ; $9$ ; $12$ ; $7$ ; $10$. Calculer le nombre moyen de pas par
    jour.
  \item Dans une classe, $25$ élèves ont indiqué le nombre d'activités auxquelles ils
    sont inscrits. Calculer le nombre moyen d'activités.
    \begin{tableaustat}
    \begin{tabular}{|l|*{4}{c|}}
    \hline
    Nombre d'activités & $0$ & $1$ & $2$ & $3$\\\hline
    Effectif & $6$ & $10$ & $6$ & $3$\\\hline
    \end{tabular}
    \end{tableaustat}
  \item Une enquête sur une autoroute a relevé le nombre de personnes dans chaque
    voiture. Calculer le nombre moyen de personnes par voiture.
    \begin{tableaustat}
    \begin{tabular}{|l|*{4}{c|}}
    \hline
    Nombre de personnes & $1$ & $2$ & $3$ & $4$\\\hline
    Fréquence & $55\,\%$ & $30\,\%$ & $10\,\%$ & $5\,\%$\\\hline
    \end{tabular}
    \end{tableaustat}
\end{enumerate}

\exo[Moyenne avec coefficients]
Inès a obtenu les notes suivantes ce trimestre.
\begin{tableaustat}
\begin{tabular}{|l|*{4}{c|}}
\hline
Note & $12$ & $8$ & $15$ & $10$\\\hline
Coefficient & $2$ & $1$ & $3$ & $2$\\\hline
\end{tabular}
\end{tableaustat}
\begin{enumerate}[label=\alph*)]
  \item Calculer sa moyenne. Pourquoi est-elle différente de la moyenne des quatre
    notes sans coefficient ?
  \item Il reste un devoir de coefficient $2$. Quelle note doit-elle obtenir pour
    avoir exactement $12{,}5$ de moyenne ?
\end{enumerate}

\exo[Moyenne de deux classes]
Au même devoir, les $32$ élèves de la seconde A ont eu $11{,}5$ de moyenne, et les
$28$ élèves de la seconde B, $13$.
\begin{enumerate}[label=\alph*)]
  \item Calculer la moyenne des $60$ élèves.
  \item Expliquer pourquoi ce n'est pas $\dfrac{11{,}5+13}{2}$.
  \item Les $30$ élèves de la seconde C ont eu $10$ de moyenne. Calculer la moyenne
    des $90$ élèves, arrondie au centième.
\end{enumerate}

\exo[Trajets des salariés]
Une entreprise a relevé le temps de trajet domicile-travail de ses $80$ salariés.
\begin{tableaustat}
\begin{tabular}{|l|*{4}{c|}}
\hline
Durée (min) & $\intfo{0}{15}$ & $\intfo{15}{30}$ & $\intfo{30}{45}$ & $\intfo{45}{60}$\\\hline
Effectif & $18$ & $30$ & $24$ & $8$\\\hline
\end{tabular}
\end{tableaustat}
\begin{enumerate}[label=\alph*)]
  \item Donner le centre de chaque classe.
  \item Estimer le temps de trajet moyen, arrondi au dixième de minute.
  \item Quelle proportion des salariés mettent moins de $30$ min ?
\end{enumerate}

\exo[Linéarité de la moyenne]
\begin{enumerate}[label=\alph*)]
  \item Le prix moyen de cinq articles est $24$ €. Chaque prix augmente de
    $5\,\%$. Quel est le nouveau prix moyen ?
  \item Le magasin fait ensuite une réduction de $2$ € sur chaque article. Quel est
    alors le prix moyen ?
  \item Une semaine, la température moyenne a été de $15$ \textdegree C. Pour passer en
    degrés Fahrenheit, on multiplie par $1{,}8$ puis on ajoute $32$. Calculer la
    température moyenne en degrés Fahrenheit, sans convertir chaque température.
\end{enumerate}

\exo[Une note de plus]
Tom a $11$ de moyenne sur quatre notes de même coefficient.
\begin{enumerate}[label=\alph*)]
  \item Calculer la somme de ses quatre notes.
  \item Il obtient $16$ au devoir suivant. Calculer sa nouvelle moyenne.
  \item Quelle note doit-il obtenir au sixième devoir pour avoir $12{,}5$ de
    moyenne ?
\end{enumerate}

\partiefiche{Étendue, médiane, quartiles}

\exo[Médiane et quartiles de deux listes]
Pour chaque série, déterminer l'étendue, la médiane, $Q_1$ et $Q_3$.
\begin{enumerate}[label=\alph*)]
  \item Série $A$, âges des $11$ joueurs d'une équipe : $19$ ; $31$ ; $24$ ; $27$ ;
    $22$ ; $29$ ; $33$ ; $25$ ; $21$ ; $28$ ; $26$.
  \item Série $B$, buts marqués lors de $10$ matchs : $2$ ; $0$ ; $3$ ; $1$ ; $1$ ;
    $4$ ; $2$ ; $0$ ; $5$ ; $2$.
  \item Traduire par une phrase le troisième quartile de la série $B$.
\end{enumerate}

\exo[Séances de cinéma]
On a demandé à $40$ personnes combien de fois elles sont allées au cinéma ce
trimestre.
\begin{tableaustat}
\begin{tabular}{|l|*{6}{c|}}
\hline
Nombre de séances & $0$ & $1$ & $2$ & $3$ & $4$ & $5$\\\hline
Effectif & $8$ & $12$ & $9$ & $6$ & $3$ & $2$\\\hline
\end{tabular}
\end{tableaustat}
\begin{enumerate}[label=\alph*)]
  \item Calculer les effectifs cumulés croissants.
  \item Déterminer la médiane, $Q_1$ et $Q_3$.
  \item Calculer l'écart interquartile.
  \item Calculer le nombre moyen de séances.
\end{enumerate}

\exo[Le salaire du patron]
Voici les salaires mensuels nets, en euros, des $8$ salariés d'une petite
entreprise : $1~800$ ; $1~900$ ; $2~000$ ; $2~100$ ; $2~100$ ; $2~200$ ; $2~300$ ;
$2~400$.
\begin{enumerate}[label=\alph*)]
  \item Calculer la moyenne, la médiane et l'étendue de ces salaires.
  \item On ajoute le salaire du patron, $9~600$ €. Calculer de nouveau la moyenne,
    la médiane et l'étendue, au centime près.
  \item Quel indicateur décrit le mieux le salaire « typique » d'un salarié ?
    Expliquer.
\end{enumerate}

\exo[Âges d'une association]
Le tableau répartit les $50$ membres d'une association selon leur âge.
\begin{tableaustat}
\begin{tabular}{|l|*{6}{c|}}
\hline
Âge & $\intfo{10}{20}$ & $\intfo{20}{30}$ & $\intfo{30}{40}$ & $\intfo{40}{50}$
  & $\intfo{50}{60}$ & $\intfo{60}{70}$\\\hline
Effectif & $6$ & $14$ & $9$ & $8$ & $7$ & $6$\\\hline
\end{tabular}
\end{tableaustat}
\begin{enumerate}[label=\alph*)]
  \item Calculer les effectifs cumulés croissants.
  \item Déterminer la classe médiane.
  \item Estimer la médiane, au dixième d'année, en supposant les âges répartis
    uniformément dans cette classe.
  \item Estimer l'âge moyen à l'aide des centres des classes.
\end{enumerate}

\exo[Contrôle d'une chaîne de remplissage]
Une machine remplit des pots de $500$ g. On a pesé $100$ pots.
\begin{tableaustat}
\begin{tabular}{|l|*{7}{c|}}
\hline
Masse (g) & $497$ & $498$ & $499$ & $500$ & $501$ & $502$ & $503$\\\hline
Effectif & $2$ & $5$ & $14$ & $40$ & $24$ & $10$ & $5$\\\hline
\end{tabular}
\end{tableaustat}
La machine est bien réglée si la moyenne est à moins de $0{,}5$ g de $500$ g et si
moins de $10\,\%$ des pots sont en dehors de $\intff{498}{502}$.
\begin{enumerate}[label=\alph*)]
  \item On retire $500$ à chaque masse. Calculer la moyenne de cette nouvelle série,
    puis en déduire la masse moyenne.
  \item Déterminer la médiane, $Q_1$ et $Q_3$.
  \item Calculer le pourcentage de pots en dehors de $\intff{498}{502}$.
  \item La machine est-elle bien réglée ?
\end{enumerate}

\partiefiche{Écart type}

\exo[Lectures régulières]
Voici le nombre de pages lues chaque jour, pendant cinq jours, par Léo et par Zoé.
\begin{center}
Léo : $6$ ; $8$ ; $10$ ; $12$ ; $14$ \qquad Zoé : $2$ ; $18$ ; $10$ ; $4$ ; $16$
\end{center}
\begin{enumerate}[label=\alph*)]
  \item Calculer à la main la moyenne et l'écart type de la série de Léo, arrondi au
    centième.
  \item Déterminer à la calculatrice la moyenne et l'écart type de la série de Zoé.
  \item Comparer les deux séries par une phrase.
\end{enumerate}

\exo[Lancers de trois pièces]
On a lancé $24$ fois trois pièces de monnaie, en notant chaque fois le nombre de
« pile ».
\begin{tableaustat}
\begin{tabular}{|l|*{4}{c|}}
\hline
Nombre de « pile » & $0$ & $1$ & $2$ & $3$\\\hline
Effectif & $3$ & $9$ & $9$ & $3$\\\hline
\end{tabular}
\end{tableaustat}
\begin{enumerate}[label=\alph*)]
  \item Calculer la moyenne de cette série.
  \item Recopier et compléter :
    $\sigma=\sqrt{\dfrac{3\times(0-1{,}5)^{2}+9\times(\ldots)^{2}+\ldots}{24}}$.
  \item Calculer l'écart type, arrondi au centième, puis le vérifier à la
    calculatrice.
\end{enumerate}

\exo[Choisir avec les indicateurs]
\begin{enumerate}[label=\alph*)]
  \item Deux machines remplissent des paquets de $250$ g. Machine $A$ :
    $\overline{x}=250{,}2$ g et $\sigma=0{,}8$ g ; machine $B$ : $\overline{x}=250{,}1$ g
    et $\sigma=3{,}1$ g. Laquelle est la plus régulière ? Répondre par une phrase.
  \item Deux secondes ont la même moyenne, $11$, à un devoir. L'écart type vaut $2$
    en seconde D et $5$ en seconde E. Dans quelle classe les notes sont-elles le plus
    dispersées ?
  \item Une série de salaires a pour médiane $2~100$ € et pour écart interquartile
    $300$ €. Écrire une phrase qui traduit chacun de ces deux nombres.
\end{enumerate}

\exo[Sélection au tir à l'arc]
Deux archers, $X$ et $Y$, ont tiré $20$ flèches chacun.
\begin{tableaustat}
\begin{tabular}{|l|*{5}{c|}}
\hline
Points & $0$ & $4$ & $6$ & $8$ & $10$\\\hline
Flèches de $X$ & $1$ & $3$ & $6$ & $7$ & $3$\\\hline
Flèches de $Y$ & $3$ & $3$ & $4$ & $4$ & $6$\\\hline
\end{tabular}
\end{tableaustat}
\begin{enumerate}[label=\alph*)]
  \item Calculer la moyenne de chaque archer.
  \item Déterminer la médiane de chaque archer.
  \item Déterminer à la calculatrice l'écart type de chaque archer, arrondi au
    centième.
  \item Quel archer est le plus régulier ? Lequel choisir pour une compétition où
    seule la régularité compte ?
\end{enumerate}

\partiefiche{Diagrammes}

\exo[Matière préférée]
Dans une classe, $32$ élèves ont dit leur matière préférée ; deux ont répondu
« ne se prononce pas ».
\begin{tableaustat}
\begin{tabular}{|l|*{5}{c|}}
\hline
Matière & maths & langues & histoire & sciences & EPS\\\hline
Effectif & $6$ & $5$ & $4$ & $7$ & $8$\\\hline
\end{tabular}
\end{tableaustat}
\begin{enumerate}[label=\alph*)]
  \item Calculer la fréquence de chaque matière, puis l'angle de son secteur.
  \item Construire le diagramme circulaire.
  \item Quel autre diagramme aurait convenu ? Pourquoi ne peut-on pas tracer
    d'histogramme ?
\end{enumerate}

\exo[Lire un diagramme circulaire]
Un club omnisports compte $240$ adhérents, répartis selon leur sport principal.
\begin{center}
\begin{tikzpicture}
  \draw[secteur,fill=coulCourbe!20] (0,0) -- (90:2) arc (90:210:2) -- cycle;
  \draw[secteur,fill=coulCourbeC!20] (0,0) -- (210:2) arc (210:270:2) -- cycle;
  \draw[secteur,fill=coulCourbeB!20] (0,0) -- (270:2) arc (270:315:2) -- cycle;
  \draw[secteur,fill=coulCourbe!8] (0,0) -- (315:2) arc (315:405:2) -- cycle;
  \draw[secteur,fill=coulCourbeC!8] (0,0) -- (45:2) arc (45:90:2) -- cycle;
  \node[font=\footnotesize,align=center] at (150:1.15) {football\\$120^{\circ}$};
  \node[font=\footnotesize,align=center] at (240:1.25) {basket\\$60^{\circ}$};
  \node[font=\footnotesize,align=center] at (292:1.35) {tennis\\$45^{\circ}$};
  \node[font=\footnotesize,align=center] at (0:1.15) {handball\\$90^{\circ}$};
  \node[font=\footnotesize,align=center] at (67:1.4) {escalade\\$45^{\circ}$};
\end{tikzpicture}
\end{center}
\begin{enumerate}[label=\alph*)]
  \item Calculer le nombre d'adhérents de chaque sport.
  \item Quel pourcentage des adhérents pratiquent le handball ?
\end{enumerate}

\exo[Histogramme et polygone]
Une entreprise a relevé le temps de trajet domicile-travail de ses $80$ salariés.
\begin{tableaustat}
\begin{tabular}{|l|*{4}{c|}}
\hline
Durée (min) & $\intfo{0}{15}$ & $\intfo{15}{30}$ & $\intfo{30}{45}$ & $\intfo{45}{60}$\\\hline
Effectif & $18$ & $30$ & $24$ & $8$\\\hline
\end{tabular}
\end{tableaustat}
\begin{enumerate}[label=\alph*)]
  \item Construire l'histogramme de cette série.
  \item Calculer les fréquences cumulées croissantes, en pourcentage.
  \item Construire le polygone des fréquences cumulées.
  \item Lire sur le polygone une estimation de la médiane.
  \item Retrouver cette estimation par le calcul, avec la classe médiane.
\end{enumerate}

\exo[Lire un histogramme]
L'histogramme répartit des lycéens selon leur temps de sport hebdomadaire, en heures
($1$ carreau pour $1$ élève).
\begin{center}
\begin{tikzpicture}[x=0.6cm,y=0.25cm]
  \quadrillage{0}{0}{11}{13}
  \foreach \a/\h in {0/5,2/12,4/9,6/6,8/3} \draw[barre] (\a,0) rectangle ({\a+2},\h);
  \draw[-{Stealth[length=2mm]},thick] (0,0)--(11.5,0) node[above left,font=\footnotesize]{h};
  \draw[-{Stealth[length=2mm]},thick] (0,0)--(0,13.6) node[left,font=\footnotesize]{effectif};
  \foreach \k in {0,2,4,6,8,10}
    \draw (\k,0.3)--(\k,-0.3) node[below,font=\scriptsize]{$\k$};
  \foreach \k in {5,10}
    \draw (0.12,\k)--(-0.12,\k) node[left,font=\scriptsize]{$\k$};
\end{tikzpicture}
\end{center}
\begin{enumerate}[label=\alph*)]
  \item Dresser le tableau des effectifs.
  \item Combien d'élèves ont été interrogés ?
  \item Quel pourcentage des élèves font au moins $6$ h de sport, au dixième près ?
  \item Estimer le temps de sport moyen à l'aide des centres des classes.
\end{enumerate}

\columnbreak
\partiefiche{Croisement de deux variables qualitatives}

\exo[Petit-déjeuner]
Une enquête porte sur les $800$ élèves d'un lycée : ont-ils pris un petit-déjeuner ce
matin ?
\begin{tableaustat}
\begin{tabular}{|l|c|c|c|c|}
\hline
 & seconde & première & terminale & Total\\\hline
Oui & $210$ & & $156$ & \\\hline
Non & $90$ & & $84$ & \\\hline
Total & $300$ & $260$ & $240$ & $800$\\\hline
\end{tabular}
\end{tableaustat}
\begin{enumerate}[label=\alph*)]
  \item Parmi les élèves de première, $65\,\%$ ont pris un petit-déjeuner. Compléter
    le tableau.
  \item Calculer la fréquence marginale des élèves qui ont pris un petit-déjeuner.
  \item Calculer la fréquence des élèves qui ont pris un petit-déjeuner parmi les
    terminales.
  \item Calculer la fréquence des terminales parmi les élèves qui n'ont pas pris de
    petit-déjeuner. Traduire le résultat par une phrase.
\end{enumerate}

\exo[Lire les phrases des médias]
Au recensement de $2017$, en France :
\begin{itemize}
  \item[(1)] $4{,}5\,\%$ des actifs qui habitent une ville-centre vont travailler à vélo ;
  \item[(2)] chez les actifs des communes hors de l'attraction des villes, la part du vélo tombe à
    $1{,}1\,\%$ ;
  \item[(3)] $73{,}7\,\%$ des actifs vont travailler en voiture.
\end{itemize}
\begin{enumerate}[label=\alph*)]
  \item Pour chaque phrase, dire si le pourcentage est une fréquence marginale ou
    conditionnelle, et donner sa population de référence.
  \item Un article affirme : « $4{,}5\,\%$ des cyclistes habitent une ville-centre ».
    Expliquer pourquoi cette phrase ne dit pas la même chose que la phrase (1).
\end{enumerate}

\exo[Le vélo et la distance]
En $2017$, $24~600$ milliers d'actifs allaient travailler à moins de $150$ km : $3~700$
milliers à $2$ km ou moins, $4~500$ milliers de $2$ à $5$ km et $16~400$ milliers à
plus de $5$ km. Le vélo était le mode principal de $5\,\%$ des premiers, de $5{,}1\,\%$
des deuxièmes et de $0{,}9\,\%$ des derniers.
\begin{enumerate}[label=\alph*)]
  \item Calculer le nombre de cyclistes, en milliers arrondis à l'unité, pour chaque
    distance, puis au total.
  \item Calculer la fréquence marginale des cyclistes parmi les actifs, au dixième
    de pourcentage.
  \item Quelle part des cyclistes habitent à $5$ km ou moins de leur travail ?
    Répondre par une phrase.
\end{enumerate}

\exo[Satisfaction sur une ligne de bus]
Une enquête interroge $400$ usagers d'une ligne de bus : $250$ sont des usagers
réguliers, les autres occasionnels. $70\,\%$ des réguliers et $80\,\%$ des
occasionnels se disent satisfaits.
\begin{enumerate}[label=\alph*)]
  \item Dresser le tableau croisé des effectifs (régulier ou occasionnel ; satisfait
    ou non).
  \item Calculer la fréquence marginale des usagers satisfaits.
  \item Quelle part des usagers non satisfaits sont des réguliers ?
  \item Un journal titre : « $70\,\%$ des usagers sont satisfaits ». Est-ce exact ?
  \item On choisit un usager au hasard. On note $R$ : « l'usager est régulier » et
    $S$ : « l'usager est satisfait ». Calculer $P(R\cap S)$ et $P_S\big(\overline{R}\big)$.
\end{enumerate}

\end{multicols}

\end{document}
